\documentclass[11pt,a4paper]{article}
\usepackage[utf8]{inputenc}
\usepackage[T1]{fontenc}
\usepackage{lmodern}
\usepackage{microtype}
\usepackage{amsmath,amssymb,amsthm}
\usepackage{graphicx}
\usepackage{booktabs}
\usepackage{longtable}
\usepackage{array}
\usepackage{multirow}
\usepackage{hyperref}
\usepackage[round]{natbib}
\usepackage{geometry}
\usepackage{fancyhdr}
\usepackage{setspace}
\usepackage{caption}
\usepackage{subcaption}

% Page setup
\geometry{margin=1in}
\onehalfspacing

% Header/Footer
\pagestyle{fancy}
\fancyhf{}
\fancyhead[L]{Climate Risk Premium Working Paper}
\fancyhead[R]{\today}
\fancyfoot[C]{\thepage}

% Title information
\title{Quantifying the Climate Risk Premium: A Case Study of the Samcheok Blue Power Plant in South Korea}
\author{Jinsu Park\\
\texttt{PLANiT Institute, Seoul, South Korea}\\
\texttt{jinsu@planit.institute}}
\date{Working Paper -- February 2026}

\begin{document}

\maketitle

\begin{abstract}
This paper presents a comprehensive financial modeling framework for quantifying the Climate Risk Premium (CRP) of coal-fired power infrastructure, using the 2.1 GW Samcheok Blue Power plant in South Korea as a case study. By integrating three independent data sources—Korea Power Supply Plan (MOTIE), CLIMADA physical hazards, and KIS credit rating methodology—we demonstrate that government policy, not physical climate change, is the primary driver of coal asset stranding in Korea. Under the enhanced 11th Basic Plan (2040 coal phase-out), the Samcheok plant's NPV drops from +\$3,098M to -\$2,999M and its credit rating collapses from AA to CC, generating a Climate Risk Premium of 1,020 basis points. Physical risks alone have negligible impact (\$3-5M NPV reduction), confirming that transition risk is the dominant stranding channel. Our model introduces a novel "credit rating death spiral" mechanism that captures the non-linear feedback loop between climate risks, cash flows, and financing costs. These findings have significant implications for investors, policymakers, and rating agencies in assessing climate-related financial risks in the energy sector.
\end{abstract}

\textbf{Keywords:} Climate risk premium, stranded assets, coal power, credit rating, transition risk, physical risk, South Korea

\textbf{JEL Classification:} G32, Q54, Q43, L94

\section{Introduction}

\subsection{The Samcheok Paradox}

South Korea faces a critical dilemma: while committing to carbon neutrality by 2050, the country recently commissioned the 2.1 GW Samcheok Blue Power plant in 2024—likely the last coal-fired power plant in its history. This contradiction presents a unique case study for analyzing "stranded asset" risk in real-time, providing an opportunity to quantify the financial implications of climate policy on energy infrastructure.

The Samcheok plant represents a \$6.1 billion investment that faces unprecedented uncertainty from multiple climate-related risk channels. Understanding how these risks translate into financial metrics is crucial for investors, policymakers, and financial institutions managing climate-related exposures.

\subsection{Research Questions}

This study addresses four key research questions:
\begin{enumerate}
\item How do government energy policies translate into plant-level financial impacts?
\item How do physical climate hazards (wildfire, flood, sea level rise) affect project economics?
\item How do climate risks trigger credit rating downgrades and financing cost increases?
\item What is the total "Climate Risk Premium" investors should demand?
\end{enumerate}

\subsection{Key Contributions}

Our research makes several novel contributions:
\begin{itemize}
\item \textbf{Integrated Risk Framework:} Combines physical and transition climate risks with financial modeling
\item \textbf{Credit Rating Death Spiral:} Introduces a non-linear feedback mechanism between climate risks and financing costs
\item \textbf{Policy-Dominance Finding:} Demonstrates that transition risk dominates physical risk for Korean coal assets
\item \textbf{Quantitative CRP Estimation:} Provides a methodology for calculating climate risk premiums in basis points
\end{itemize}

\section{Literature Review}

\subsection{Climate Risk in Energy Finance}

The financial implications of climate change have become increasingly important in energy finance research \citep{caldecott2021stranded}. Studies have identified two primary risk channels: physical risks from climate-related events and transition risks from policy and technological changes \citep{daumas2024financial}.

\subsection{Stranded Assets Research}

Stranded asset risk has been extensively studied in the context of fossil fuel infrastructure \citep{bressan2024asset}. However, most research focuses on aggregate portfolio-level analysis rather than plant-specific financial modeling \citep{hoffart2024energy}.

\subsection{Credit Risk and Climate}

Recent research has begun exploring the relationship between climate risks and credit ratings \citep{jung2021crisk}. However, existing models often fail to capture the dynamic feedback loops between climate impacts and financing costs \citep{albanese2024effects}.

\section{Methodology}

\subsection{Model Architecture}

Our integrated framework combines three modules:
\begin{enumerate}
\item \textbf{Physical Risk Module:} Quantifies operational impacts from climate hazards
\item \textbf{Transition Risk Module:} Models policy-driven changes in dispatch and carbon costs
\item \textbf{Financial Impact Module:} Calculates NPV, IRR, and credit rating implications
\end{enumerate}

\subsection{Data Sources}

\subsubsection{Korea Power Supply Plan (MOTIE)}
Official government coal dispatch trajectories for 2024-2050, including capacity factors and phase-out schedules under different policy scenarios.

\subsubsection{CLIMADA Physical Hazards}
Spatially-explicit wildfire, flood, and sea level rise data at 4.5 km resolution, calibrated for the Samcheok location (37.44°N, 129.17°E).

\subsubsection{KIS Credit Rating Methodology}
Korean credit rating agency quantitative grid for Independent Power Producers (IPPs), including capacity, coverage, and leverage metrics.

\subsection{Financial Model}

The core financial model calculates project cash flows using standard project finance methodology:

\begin{equation}
\text{NPV} = \sum_{t=1}^{T} \frac{\text{CF}_t}{(1 + \text{WACC})^t} - I_0
\end{equation}

where:
\begin{align}
\text{CF}_t &= \text{EBIT}_t \times (1 - \tau) + \text{Depreciation}_t - \text{Capex}_t - \Delta\text{WC}_t \\
\tau &= 24\% \text{ (Korean corporate tax rate)} \\
\text{WACC} &= \left(\frac{E}{V} \times r_e\right) + \left(\frac{D}{V} \times r_d \times (1 - \tau)\right)
\end{align}

\subsection{Climate Risk Premium Calculation}

The Climate Risk Premium (CRP) is defined as:

\begin{equation}
\text{CRP} = \text{Spread}(R_{\text{risk}}) - \text{Spread}(R_{\text{baseline}}) + \text{Expected Loss Spread}
\end{equation}

where $R = f(\text{DSCR}, \text{EBITDA/Interest}, \text{Net Debt/EBITDA}, \dots)$ represents the credit rating function.

\section{Results}

\subsection{Scenario Analysis}

Table \ref{tab:scenarios} presents the results across 11 scenarios, including baseline, physical risk, transition risk, and combined scenarios.

\begin{table}[htbp]
\centering
\caption{Scenario Analysis Results}
\label{tab:scenarios}
\begin{tabular}{lrrrrr}
\toprule
\textbf{Scenario} & \textbf{NPV (\$M)} & \textbf{IRR} & \textbf{Min DSCR} & \textbf{Rating} & \textbf{CRP (bps)} \\
\midrule
Baseline & 3,098 & 11.99\% & 1.86× & AA & -50 \\
Moderate Transition & 2,034 & 10.55\% & 1.65× & A & 0 \\
Aggressive Transition & -75 & 7.04\% & 1.33× & A & 0 \\
Moderate Physical & 3,095 & 11.99\% & 1.85× & AA & -50 \\
High Physical & 3,093 & 11.99\% & 1.85× & AA & -50 \\
Combined Moderate & 2,031 & 10.55\% & 1.65× & A & 0 \\
Combined Aggressive & -78 & 7.03\% & 1.33× & A & 0 \\
Low Demand & 494 & 8.21\% & 1.17× & BBB & 85 \\
Severe Drought & 3,099 & 11.99\% & 1.86× & AA & -50 \\
Enhanced 11th Plan & -2,999 & -8.09\% & 0.07× & CC & 1,020 \\
Enhanced Combined & -3,000 & -8.09\% & 0.07× & CC & 1,020 \\
\bottomrule
\end{tabular}
\end{table}

\subsection{Key Findings}

\subsubsection{Policy Dominates Physical Risk}
The Enhanced 11th Plan (2040 coal phase-out) destroys \$6.1B in value (NPV swing from +\$3,098M to -\$2,999M), while physical risks have negligible financial impact (\$3-5M NPV reduction).

\subsubsection{Credit Rating Collapse}
Baseline AA rating drops to CC under the 2040 phase-out scenario, with DSCR collapsing from 1.86x to 0.07x.

\subsubsection{Climate Risk Premium}
1,020 bps under the most severe policy scenario, reflecting a 10.2\% additional cost of capital.

\subsubsection{Physical Risk Minimal for Korea}
Korea-specific wildfire (0.055\%) and flood (0.003\%) risks are orders of magnitude smaller than transition risk, consistent with Korea's temperate geography.

\subsection{Credit Rating Death Spiral}

Figure \ref{fig:death_spiral} illustrates the non-linear feedback loop:

\begin{enumerate}
\item Climate risks reduce revenue and cash flows
\item Lower cash flows reduce Debt Service Coverage Ratio (DSCR)
\item Lower DSCR triggers credit rating downgrades
\item Lower ratings increase cost of debt (spread widens)
\item Higher interest expense further reduces cash flows
\item Loop repeats until technical default
\end{enumerate}

\section{Discussion}

\subsection{Policy Implications}

Our findings have significant implications for energy policy design:

\begin{itemize}
\item \textbf{Stranded Asset Risk is Real:} The Enhanced 11th Basic Plan creates a \$6.1 billion NPV swing for the Samcheok plant
\item \textbf{Early Retirement is Optimal:} Under the 2040 phase-out, DSCR collapses to 0.07x by mid-life
\item \textbf{Just Transition Finance Needed:} The 1,020 bps CRP renders new coal unfinanceable
\end{itemize}

\subsection{Investment Implications}

For investors and financial institutions:

\begin{itemize}
\item \textbf{Climate Risk Premium:} 1,020 bps under severe policy scenarios must be priced into new investments
\item \textbf{Rating Agency Adaptation:} Current static methodologies fail to capture forward-looking policy risk
\item \textbf{Portfolio Management:} Transition risk dominates physical risk for Korean coal assets
\end{itemize}

\subsection{Methodological Contributions}

Our study introduces several methodological innovations:

\begin{itemize}
\item \textbf{Integrated Risk Assessment:} Combines physical and transition risks in a unified framework
\item \textbf{Dynamic Credit Modeling:} Captures feedback loops between climate impacts and financing costs
\item \textbf{Plant-Level Analysis:} Provides detailed financial modeling rather than aggregate portfolio analysis
\end{itemize}

\section{Conclusion}

This study demonstrates that government policy, not physical climate change, is the primary driver of coal asset stranding in South Korea. The Samcheok Blue Power plant case reveals a potential \$6.1 billion value destruction under the Enhanced 11th Basic Plan, with credit ratings collapsing from AA to CC and financing costs increasing by over 1,000 basis points.

Our "credit rating death spiral" mechanism provides a new framework for understanding how climate risks propagate through financial systems. The finding that transition risk dominates physical risk has important implications for climate risk assessment, investment decisions, and policy design.

Future research should extend this framework to other asset classes, geographic regions, and policy contexts to develop a more comprehensive understanding of climate-related financial risks.

\section*{Acknowledgments}

The author thanks the PLANiT Institute for research support, Solutions for Our Climate (SFOC) for Korean coal policy data, the ETH Zurich CLIMADA Team for open-source hazard modeling tools, and Korea Investors Service (KIS) for credit rating methodology guidance.

\bibliographystyle{apalike}
\bibliography{references}

\appendix

\section{Model Parameters}

\subsection{Plant Specifications}
\begin{itemize}
\item Capacity: 2,100 MW
\item Location: 37.44°N, 129.17°E (Gangwon Province, South Korea)
\item Investment Cost: \$6.1 billion
\item Commissioning: 2024
\end{itemize}

\subsection{Financial Assumptions}
\begin{itemize}
\item Debt/Equity Ratio: 70/30
\item Corporate Tax Rate: 24\%
\item WACC: 6.75\% (baseline)
\end{itemize}

\subsection{Physical Risk Parameters}
\begin{itemize}
\item Wildfire Base Rate: 0.055\% (from Kim et al. 2025)
\item Flood Base Rate: 0.003\% (from Kang \& Lee 2024)
\item Sea Level Rise Derate: 0.22\%/meter (from Van Vliet et al. 2016)
\end{itemize}

\end{document}